                         IMC2010, Blagoevgrad, Bulgaria
                                        Day 2, July 27, 2010

Problem 1. (a) A sequence x1 , x2 , . . . of real numbers satisfies

                                     xn+1 = xn cos xn        for all n ≥ 1 .

Does it follow that this sequence converges for all initial values x1 ?
   (b) A sequence y1 , y2, . . . of real numbers satisfies

                                      yn+1 = yn sin yn       for all n ≥ 1 .

Does it follow that this sequence converges for all initial values y1 ?
Solution 1. (a) NO. For example, for x1 = π we have xn = (−1)n−1 π, and the sequence is divergent.
    (b) YES. Notice that |yn | is nonincreasing and hence converges to some number a ≥ 0.
    If a = 0, then lim yn = 0 and we are done. If a > 0, then a = lim |yn+1| = lim |yn sin yn | = a·| sin a|,
so sin a = ±1 and a = (k + 12 )π for some nonnegative integer k.
    Since the sequence |yn | is nonincreasing, there exists an index n0 such that (k + 21 )π ≤ |yn | <
(k
 + 1)π    for all n > n0 . Then all the numbers yn0 +1 , yn0 +2 , . . . lie in the union of the intervals
        1
 (k + 2 )π, (k + 1)π and − (k + 1)π, −(k + 12 )π .
    Depending on the parity of k, in one of the intervals (k+ 21 )π, (k+1)π and −(k+1)π, −(k+ 21 )π
                                                                                                               

the values of the sine function is positive; denote this interval by I+ . In the other interval the sine
function is negative; denote this interval by I− . If yn ∈ I− for some n > n0 then yn and yn+1 = yn sin yn
have opposite signs, so yn+1 ∈ I+ . On the other hand, if If yn ∈ I+ for some n > n0 then yn and yn+1
have the same sign, so yn+1 ∈ I+ . In both cases, yn+1 ∈ I+ .
    We obtained that the numbers yn0 +2 , yn0 +3 , . . . lie in I+ , so they have the same sign. Since |yn | is
convergent, this implies that the sequence (yn ) is convergent as well.
Solution 2 for part (b). Similarly to the first solution, |yn | → a for some real number a.
    Notice that t · sin t = (−t) sin(−t) = |t| sin |t| for all real t, hence yn+1 = |yn | sin |yn | for all n ≥ 2.
Since the function t 7→ t sin t is continuous, yn+1 = |yn | sin |yn | → |a| sin |a| = a.
Problem 2. Let a0 , a1 , . . . , an be positive real numbers such that ak+1 − ak ≥ 1 for all k =
0, 1, . . . , n − 1. Prove that
                                                                           
                     1          1                1           1        1           1
                1+       1+            ··· 1+          ≤ 1+       1+      ··· 1+     .
                     a0      a1 − a0          an − a0       a0       a1          an

Solution. Apply induction on n. Considering the empty product as 1, we have equality for n = 0.
   Now assume that the statement is true for some n and prove it for n + 1. For n + 1, the statement
can be written as the sum of the inequalities
                                                                           
                     1         1                  1              1              1
                1+       1+           ··· 1+            ≤ 1+         ··· 1+
                    a0      a1 − a0            an − a0          a0             an
(which is the induction hypothesis) and
                                                                       
          1          1                 1            1          1           1       1
              1+            ··· 1+            ·           ≤ 1+      ··· 1+      ·      .                      (1)
         a0       a1 − a0           an − a0     an+1 − a0      a0          an     an+1
Hence, to complete the solution it is sufficient to prove (1).

                                                         1
   To prove (1), apply a second induction. For n = 0, we have to verify
                                                         
                                   1      1             1 1
                                      ·        ≤ 1+           .
                                   a0 a1 − a0          a0 a1

Multiplying by a0 a1 (a1 − a0 ), this is equivalent with

                                         a1 ≤ (a0 + 1)(a1 − a0 )
                                             a0 ≤ a0 a1 − a20
                                              1 ≤ a1 − a0 .

   For the induction step it is sufficient that
                                                               
                                    1         an+1 − a0       1       an+1
                          1+                ·           ≤ 1+        ·      .
                               an+1 − a0      an+2 − a0      an+1     an+2

Multiplying by (an+2 − a0 )an+2 ,

                             (an+1 − a0 + 1)an+2 ≤ (an+1 + 1)(an+2 − a0 )
                                        a0 ≤ a0 an+2 − a0 an+1
                                           1 ≤ an+2 − an+1 .

Remark 1. It is easy to check from the solution that equality holds if and only if ak+1 − ak = 1 for
all k.
Remark 2. The statement of the problem is a direct corollary of the identity
                            n                        ! Y  n          
                           X     1 Y             1                   1
                       1+                  1+           =       1+       .
                           i=0
                                 xi j6=i      xj − xi      i=0
                                                                     xi

Problem 3. Denote by Sn the group of permutations of the sequence (1, 2, . . . , n). Suppose that
G is a subgroup of Sn , such that for every π ∈ G \ {e} there exists a unique k ∈ {1, 2, . . . , n} for
which π(k) = k. (Here e is the unit element in the group Sn .) Show that this k is the same for all
π ∈ G \ {e}.
Solution. Let us consider the action of G on the set X = {1, . . . , n}. Let

                        Gx = { g ∈ G : g(x) = x} and Gx = { g(x) : g ∈ G}
be the stabilizer and the orbit of x ∈ X under this action, respectively. The condition of the problem
states that                                        [
                                             G=       Gx                                            (1)
                                                    x∈X

and
                                    Gx ∩ Gy = {e} for all x 6= y.                                  (2)
We need to prove that Gx = G for some x ∈ X.
   Let Gx1 , . . . , Gxk be the distinct orbits of the action of G. Then one can write (1) as
                                                 k
                                                 [ [
                                            G=               Gy .                                  (3)
                                                 i=1 y∈Gxi


                                                    2
   It is well known that
                                                                      |G|
                                                          |Gx| =           .                                  (4)
                                                                     |Gx |
   Also note that if y ∈ Gx then Gy = Gx and thus |Gy| = |Gx|. Therefore,
                                          |G|    |G|
                               |Gx | =        =      = |Gy | for all y ∈ Gx.                                  (5)
                                         |Gx|   |Gy|
   Combining (3), (2), (4) and (5) we get
                                                  k                            k
                                                  [ [                          X |G|
                    |G| − 1 = |G \ {e}| =                        Gy \ {e} =                 (|Gxi | − 1),
                                                  i=1 y∈Gxi                    i=1
                                                                                   |Gxi |

hence
                                                    k            
                                              1    X         1
                                          1−     =      1−          .                                         (6)
                                             |G|   i=1
                                                           |Gxi |
   If for some i, j ∈ {1, . . . , k} |Gxi |, |Gxi | ≥ 2 then
                         k                                                
                         X             1                     1            1                      1
                            1−                    ≥       1−         + 1−          =1>1−
                         i=1
                                     |Gxi |                  2            2                     |G|

which contradicts with (6), thus we can assume that

                                              |Gx1 | = . . . = |Gxk−1 | = 1.

Then from (6) we get |Gxk | = |G|, hence Gxk = G.
Problem 4. Let A be a symmetric m × m matrix over the two-element field all of whose diagonal
entries are zero. Prove that for every positive integer n each column of the matrix An has a zero
entry.
Solution. Denote by ek (1 ≤ k ≤ m) the m-dimensional vector over F2 , whose k-th entry is 1 and all
the other elements are 0. Furthermore, let u be the vector whose all entries are 1. The k-th column
of An is An ek . So the statement can be written as An ek 6= u for all 1 ≤ k ≤ m and all n ≥ 1.
    For every pair of vectors x = (x1 , . . . , xm ) and y = (y1 , . . . , ym ), define the bilinear form (x, y) =
xT y = x1 y1 + ... + xm ym . The product (x, y) has all basic properties of scalar products (except the
property that (x, x) = 0 implies x = 0). Moreover, we have (x, x) = (x, u) for every vector x ∈ F2m .
    It is also easy to check that (w, Aw) = w T Aw = 0 for all vectors w, since A is symmetric and its
diagonal elements are 0.
Lemma. Suppose that v ∈ F2m a vector such that An v = u for some n ≥ 1. Then (v, v) = 0.
Proof. Apply induction on n. For odd values of n we prove the lemma directly. Let n = 2k + 1 and
w = Ak v. Then

            (v, v) = (v, u) = (v, An v) = v T An v = v T A2k+1 v = (Ak v, Ak+1 v) = (w, Aw) = 0.

   Now suppose that n is even, n = 2k, and the lemma is true for all smaller values of n. Let
w = Ak v; then Ak w = An v = u and thus we have (w, w) = 0 by the induction hypothesis. Hence,

             (v, v) = (v, u) = v T An v = v T A2k v = (Ak v)T (Ak v) = (Ak v, Ak v) = (w, w) = 0.

The lemma is proved.

                                                                 3
   Now suppose that An ek = u for some 1 ≤ k ≤ m and positive integer n. By the Lemma, we
should have (ek , ek ) = 0. But this is impossible because (ek , ek ) = 1 6= 0.
Problem 5. Suppose that for a function f : R → R and real numbers a < b one has f (x) = 0 for
all x ∈ (a, b). Prove that f (x) = 0 for all x ∈ R if
                                             p−1       
                                             X        k
                                                 f y+     =0
                                             k=0
                                                      p

for every prime number p and every real number y.
Solution. Let N > 1 be some integer to be defined later, and consider set of real polynomials
                                                                         n              
                                              n
                                                                          X        k
              JN =       c0 + c1 x + . . . + cn x ∈ R[x]        ∀x ∈ R     ck f x +    =0 .
                                                                       k=0
                                                                                    N

Notice that 0 ∈ JN , any linear combinations of any elements in JN is in JN , and for every P (x) ∈ JN
we have xP (x) ∈ JN . Hence, JN is an ideal of the ring R[x].
                                                                            xN − 1
    By the problem’s conditions, for every prime divisors of N we have N/p           ∈ JN . Since R[x]
                                                                           x    −1
is a principal ideal domain (due to the Euclidean algorithm), the greatest common divisor of these
                                                                               xN − 1
polynomials is an element of JN . The complex roots of the polynomial N/p               are those Nth
                                                                              x    −1
roots of unity whose order does not divide N/p. The roots of the greatest common divisor is the
intersection of such sets; it can be seen that the intersection consist of the primitive Nth roots of
unity. Therefore,                      N                 
                                         x −1
                                  gcd                p N = ΦN (x)
                                        xN/p − 1
is the Nth cyclotomic polynomial. So ΦN ∈ JN , which polynomial has degree ϕ(N).
                                      ϕ(N)                                            ϕ(N)
   Now choose N in such a way that            < b − a. It is well-known that lim inf       = 0, so there
                                        N                                     N →∞     N
                                                                    ϕ(N )
exists such a value for N. Let ΦN (x) = a0 + a1 x + . . . + aϕ(N ) x      where aϕ(N ) = 1 and |a0 | = 1.
                                        ϕ(N )
                                              ak f x + Nk = 0 for all x ∈ R.
                                         P               
Then, by the definition of JN , we have
                                             k=0

   If x ∈ [b, b + N1 ), then
                                                  ϕ(N )−1
                                                   X
                                                            ak f x − ϕ(NN)−k .
                                                                            
                                     f (x) = −
                                                   k=0

On the right-hand side, all numbers x − ϕ(NN)−k lie in (a, b). Therefore the right-hand side is zero,
and f (x) = 0 for all x ∈ [b, b + N1 ). It can be obtained similarly that f (x) = 0 for all x ∈ (a − N1 , a]
as well. Hence, f = 0 in the interval (a − N1 , b + N1 ). Continuing in this fashion we see that f must
vanish everywhere.




                                                            4
